\documentclass[11pt]{article}

\usepackage[a4paper,landscape,margin=0.6in]{geometry}
\usepackage{amsmath,amssymb}
\usepackage{array}
\usepackage{booktabs}
\usepackage{longtable}
\usepackage[strings]{underscore}
\usepackage{hyperref}

% Keep only argument-wrapping or formatting helpers; \meta marks syntax placeholders, while mathematical symbols are written directly for compact, LLM-readable source.
\newcommand{\code}[1]{\texttt{#1}}
\newcommand{\meta}[1]{\ensuremath{\langle #1\rangle}}
\newenvironment{compacttable}
  {\begingroup\footnotesize\setlength{\tabcolsep}{3pt}}
  {\endgroup}

\title{Runir Module-Program Syntax and Semantics (Indexical Concept Policies)}
\author{Dominik Drexler\\\href{mailto:dominik.drexler@liu.se}{dominik.drexler@liu.se}}
\date{}

\usepackage{xcolor}
\newcommand{\alert}[1]{\textcolor{red}{#1}}
\newcommand{\unclear}[1]{\textcolor{magenta}{#1}}

\begin{document}
\maketitle

\section{Used Symbols}

\begin{compacttable}
\begin{longtable}{>{\raggedright\arraybackslash}p{0.16\textwidth}
                  >{\raggedright\arraybackslash}p{0.78\textwidth}}
\toprule
Symbol & Meaning \\
\midrule
\endhead
\meta{X} & Runir syntax placeholder for the serialization of semantic entity $X$; monospaced text is literal syntax \\
$P,\mathcal Q$ & one grounded planning task and a generalized-planning class of tasks \\
$M_P=\langle S_P,A_P,T_P,s_0,S_G\rangle$ & finite deterministic state model induced by task $P$ \\
$S_P,S_G$ & reachable planning states and the subset of goal states \\
$A_P,T_P$ & ground actions and transition relation $T_P\subseteq S_P\times A_P\times S_P$ \\
$s,s',s_0,a$ & concrete states, initial state, and ground action, as determined by context \\
$s\xrightarrow{a}s'$ & concrete transition from $s$ to $s'$ under action $a$ \\
$\Pi=\langle m_0,\mathcal M\rangle$ & module program with $\mathcal M=\{M\}$ and entry symbol $m_0=m$ \\
$M=\langle m,A,R,q_0,Q,\Phi,\mathcal R\rangle$ & sole module with symbol, reserved empty list $A=\emptyset$, registers, entry memory, memories, features, and rules \\
$m,q,q',q_0$ & module symbol and memory-state symbols, as determined by context \\
$A,R,Q,\Phi,\mathcal R$ & reserved empty list, category-local ordered register lists, memory-state set, local-feature set, and rule set \\
$\mathbb{B},\mathbb{N}$ & boolean truth values and natural numbers; numerical values may additionally be $\infty$ \\
$B,N$ & boolean and numerical feature expressions \\
$f,f_B,f_N,f_C,f_R$ & local feature symbols, optionally distinguished by boolean, numerical, concept, or role category \\
$r,o$ & concept and role register symbols; $r$ instead denotes a rule symbol when used as a rule \\
$\rho,\rho_C,\rho_R$ & register valuation and its concept and role components \\
$\xi=(m,q,\rho,s)$ & execution configuration with module, memory, register valuation, and planning state \\
$C,E,C_e,R_e,C_{\times},E_{\times}$ & condition list, effect list, concept expression, role expression, extended concept, and extended effect \\
$x_i,d,(d_1,d_2),\mathcal{I}$ & action argument feature symbol, selected object or pair, and relational interpretation used by feature expressions \\
$i$ & finite positional index \\
$n_r,n_q,n_F,n_R$ & finite numbers of registers, memory states, feature declarations, and rules, respectively \\
$n_{\mathrm{body}},n_x$ & finite numbers of alternative rule bodies and action arguments, respectively \\
\bottomrule
\end{longtable}
\end{compacttable}

\noindent\textbf{Top goal: Generality}---learn a module program that proves on
the scaled training frontier, executes beyond it, and terminates structurally
rather than merely fitting current tasks;
\textbf{Soft goal 1: Total syntactic feature complexity}---among equally
general correct programs, minimize the sum of the DL constructor-node counts
of all declared feature expressions;
\textbf{Soft goal 2: Memory-state control}---with equal
secondary priority, maximize control flow through memory states so
Boolean/numerical SCC projections stay small and the complete structural
termination proof becomes maximally efficient.

%\noindent\textbf{Soft goal 3: Interpretable termination}---with equal secondary
%priority, prefer rules with Do bodies over rules with Sketch bodies because
%explicit Do bodies make structural termination proofs easier to interpret.

\section{State Models}

A planning task $P$ induces a finite deterministic state model
\[
  M_P=\langle S_P,A_P,T_P,s_0,S_G\rangle,
\]
where $S_P$ is the set of reachable states, $A_P$ is the set
of ground actions, $T_P\subseteq S_P\times A_P\times
S_P$ is the transition relation, $s_0\in S_P$ is the
initial state, and $S_G\subseteq S_P$ is the set of goal states. We
write $s\xrightarrow{a}s'$ for $(s,a,s')\in T_P$.

A generalized planning problem is a class $\mathcal Q$ of tasks with a common
vocabulary and action schemas. A Runir module program denotes one finite-state
controller that must solve every solvable task in $\mathcal Q$.

\section{Module Programs}

A module program is a tuple
\[
  \Pi=\langle m_0,\mathcal M\rangle,
\]
with exactly one module: $\mathcal M=\{M\}$ and $m_0=m$, where $m$ is the
symbol of $M$.

A module is a tuple
\[
  M=\langle m,A,R,q_0,Q,\Phi,\mathcal R\rangle,
\]
where $m$ is its symbol, $A=\emptyset$ is a reserved empty list, $R$ comprises
ordered concept- and role-register lists, $q_0\in Q$ is the entry memory state,
$Q$ is a finite set of memory states, $\Phi$ is a finite set of local features,
and $\mathcal R$ is a finite set of rules. Register, memory, feature, and rule
symbols are local to the module.

\begin{compacttable}
\begin{longtable}{>{\raggedright\arraybackslash}p{0.20\textwidth}
                  >{\raggedright\arraybackslash}p{0.36\textwidth}
                  >{\raggedright\arraybackslash}p{0.38\textwidth}}
\toprule
Form & Runir Syntax Schema & Meaning \\
\midrule
\endhead
Program & \code{(:program (:entry \meta{m}) \meta{M})} & Declares the sole module as the entry module. \\
  Module & \code{(:module (:symbol \meta{m}) (:arguments) (:registers \meta{r_1} \(\ldots\) \meta{r_{n_r}}) (:entry \meta{q}) (:memory \meta{q_1} \(\ldots\) \meta{q_{n_q}}) (:features \meta{F_1} \(\ldots\) \meta{F_{n_F}}) (:rules \meta{R_1} \(\ldots\) \meta{R_{n_R}}) (:reset-spo \meta{P_1} \ldots \meta{P_{n_{spo}}}))} & Declares the sole finite-state controller. \\
Concept register & \code{(:concept \meta{r})} & Declares a local concept register assignable by a matching load body. \\
Role register & \code{(:role \meta{o})} & Declares a local role register assignable by a matching load body. \\
Memory & \code{(:memory \meta{q_1} \(\ldots\) \meta{q_{n_q}})} & Declares the finite control states of the module. \\
Boolean feature & \code{(:boolean (:symbol \meta{b}) (:expression \meta{B}))} & Declares boolean feature symbol \(b\) with expression \(B\). \\
Numerical feature & \code{(:numerical (:symbol \meta{n}) (:expression \meta{N}))} & Declares numerical feature symbol \(n\) with expression \(N\). \\
Concept feature & \code{(:concept (:symbol \meta{c}) (:expression \meta{C}))} & Declares concept feature symbol \(c\) with expression \(C\). \\
Role feature & \code{(:role (:symbol \meta{f_R}) (:expression \meta{R_e}))} & Declares role feature symbol \(f_R\) with expression \(R_e\). \\
Rule & \code{(:rule (:symbol \meta{r}) (:expression (:source-memory \meta{q}) (:target-memory \meta{q'}) \meta{\mathit{body}_1} \(\ldots\) \meta{\mathit{body}_{n_{\mathrm{body}}}}))} & Declares a control transition with one or more alternative bodies. \\
Precedence & \code{(:prec \meta{f_{C}} \meta{f_{C'}})} & Declare the reset precedence relation $\meta{f_{C}}\prec\meta{f_{C'}}$ between concepts.\\
\bottomrule
\end{longtable}
\end{compacttable}

All module sections are mandatory and occur exactly in the displayed order;
\code{(:arguments)} is always empty. Only the fields shown above are semantic
fields.
{The single-module restriction makes only this module-level section empty.}
%a \unclear{\code{:do} body retains its \code{:arguments} list of action-argument features.}
Registers, features, and rules may be empty; memory is nonempty.
{A program contains exactly one module,} and every \code{:rule} expression contains
at least one body. Feature symbols are globally unique across feature
categories within a module, and rule symbols are unique. Register symbols are
unique within each category-local namespace; memory symbols are unique within
the module. Identifiers
match \code{[A-Za-z_][A-Za-z0-9_-]*}. Line comments start with \code{;}
outside quoted strings and continue to the end of the physical line.

Each body inside one \code{:rule} is an alternative transition with the same
source and target memories.

Each module may declare at most four concept registers and at most four role
registers; the two limits are independent.

\section{Feature Valuations}

For each task $P$, module $M$, concrete state $s$, and register valuation
$\rho$, each feature $f\in\Phi$ has a value $f_P(s,\rho)$. Concept, role,
boolean, and numerical expressions are defined in
\texttt{feature-syntax-semantics.tex}. Inline feature expressions may occur
only in feature declarations. Loads, conditions, effects, and action arguments
refer to declared feature symbols.

Boolean features have values in $\mathbb{B}$. Numerical features have values in
$\mathbb{N}\cup\{\infty\}$. Concept and role features denote complete sets of
objects and object pairs.
%In a \code{:do} body, a concept feature constrains an
%action parameter by membership: every matching ground action whose object lies
%in the listed denotation is a possible transition.

\section{Conditions}

A condition is evaluated in the current valuation at $s$.
Every condition references a declared local feature symbol.
An empty \code{(:conditions)} section is true.
%Features omitted from
%\code{:effects} are unconstrained; omission does not mean \code{unchanged}.
%An empty effects section imposes no feature-value constraint (the special
%memory-only meaning of an empty Sketch body is defined below).

\begin{compacttable}
\begin{longtable}{>{\raggedright\arraybackslash}p{0.18\textwidth}
                  >{\raggedright\arraybackslash}p{0.30\textwidth}
                  >{\raggedright\arraybackslash}p{0.46\textwidth}}
\toprule
Condition & Runir Syntax Schema & Semantics \\
\midrule
\endhead
Positive boolean & \code{(positive \meta{b})} & $b_P(s,\rho)=\mathrm{true}$. \\
Negative boolean & \code{(negative \meta{b})} & $b_P(s,\rho)=\mathrm{false}$. \\
Greater than zero & \code{(greater_zero \meta{n})} & $n_P(s,\rho)>0$. \\
Equal to zero & \code{(equal_zero \meta{n})} & $n_P(s,\rho)=0$. \\
\bottomrule
\end{longtable}
\end{compacttable}

%\begin{compacttable}
%\begin{longtable}{>{\raggedright\arraybackslash}p{0.18\textwidth}
%                  >{\raggedright\arraybackslash}p{0.30\textwidth}
%                  >{\raggedright\arraybackslash}p{0.46\textwidth}}
%\toprule
%Effect & Runir Syntax Schema & Semantics on $s,s'$ \\
%\midrule
%\endhead
%Positive boolean & \code{(positive \meta{b})} & $b_P(s',\rho)=\mathrm{true}$. \\
%Negative boolean & \code{(negative \meta{b})} & $b_P(s',\rho)=\mathrm{false}$. \\
%Unchanged boolean & \code{(unchanged \meta{b})} & $b_P(s',\rho)=b_P(s,\rho)$. \\
%Increases numerical & \code{(increases \meta{n})} & $n_P(s',\rho)>n_P(s,\rho)$. \\
%Decreases numerical & \code{(decreases \meta{n})} & $n_P(s',\rho)<n_P(s,\rho)$. \\
%Unchanged numerical & \code{(unchanged \meta{n})} & $n_P(s',\rho)=n_P(s,\rho)$. \\
%\bottomrule
%\end{longtable}
%\end{compacttable}

\section{Extended Conditions and Effects}

Extended conditions check whether an action argument belongs or doesn't belong to a concept,
or whther the object stored in a register belongs or doesn't belong to a concept.
Extended effects check whether an action argument or the object stored at a register enters
or exits a concept.
An object enters (resp.\ exits) a concept across a transition $(s,s')$ if the object
doesn't (resp.\ does) belong to the denotation of the concept at $s$, and it does (resp.\ doesn't)
belongs to the denotation of the concept at $s'$.

\begin{compacttable}
\begin{longtable}{>{\raggedright\arraybackslash}p{0.18\textwidth}
                  >{\raggedright\arraybackslash}p{0.40\textwidth}
                  >{\raggedright\arraybackslash}p{0.36\textwidth}}
\toprule
Extended Condition & Runir Syntax Schema & Semantics \\
\midrule
\endhead
Belongs & \code{(belongs (:argument \meta{x_i}) (:concept \meta{f_C}))} & ${x_i} \in (f_C)_P(s,\rho)$. \\
Register belongs & \code{(belongs (:register (:concept \meta{r})) (:concept \meta{f_C}))} & $\rho_C(r) \in (f_C)_P(s,\rho)$. \\
Not belongs & \code{(not-belongs (:argument \meta{x_i}) (:concept \meta{f_C}))} & ${x_i} \notin (f_C)_P(s,\rho)$. \\
Not register belongs & \code{(not-belongs (:register (:concept \meta{r})) (:concept \meta{f_C}))} & $\rho_C(r) \notin (f_C)_P(s,\rho)$. \\
\bottomrule
\end{longtable}
\end{compacttable}

\begin{compacttable}
\begin{longtable}{>{\raggedright\arraybackslash}p{0.18\textwidth}
                  >{\raggedright\arraybackslash}p{0.40\textwidth}
                  >{\raggedright\arraybackslash}p{0.36\textwidth}}
\toprule
Extended Effect & Runir Syntax Schema & Semantics on $s,s'$ \\
\midrule
\endhead
Enter & \code{(enter (:argument \meta{x_i}) (:concept \meta{f_C}))} & $\meta{x_i} \notin (f_C)_P(s,\rho) \land \meta{x_i} \in (f_C)_P(s',\rho)$. \\
Register enter & \code{(enter (:register (:concept \meta{r})) (:concept \meta{f_C}))} & $\rho_C(r) \notin (f_C)_P(s,\rho) \land \rho_C(r) \in (f_C)_P(s',\rho)$. \\
Exit & \code{(exit (:argument \meta{x_i}) (:concept \meta{f_C}))} & $\meta{x_i} \in (f_C)_P(s,\rho) \land \meta{x_i} \notin (f_C)_P(s',\rho)$. \\
Register exit & \code{(exit (:register (:concept \meta{r})) (:concept \meta{f_C}))} & $\rho_C(r) \in (f_C)_P(s,\rho) \land \rho_C(r) \notin (f_C)_P(s',\rho)$. \\
\bottomrule
\end{longtable}
\end{compacttable}


\section{Rule Bodies}

Let $r$ be a rule with source memory $q$ and target memory $q'$. A rule body is
applicable in configuration $\xi$ only if all of its conditions and extended conditions hold,
and no object re-enters the same concept twice (checked using the concept histories
in the configuration $\xi$.
Bodies inside one rule are alternatives. %Load, Do, and Sketch successors use $q'$.

\begin{compacttable}
\begin{longtable}{>{\raggedright\arraybackslash}p{0.16\textwidth}
                  >{\raggedright\arraybackslash}p{0.36\textwidth}
                  >{\raggedright\arraybackslash}p{0.42\textwidth}}
\toprule
Body & Runir Syntax Schema & Semantics \\
\midrule
\endhead
Concept & \code{(:crule (:action "\meta{a}") (:arguments \meta{x_1} \(\ldots\) \meta{x_{n_x}}) (:conditions \meta{C}) (:xconditions \meta{C_{\times}}) (:xeffects \meta{E_{\times}})))} & Branches over every transition $s\xrightarrow{a(o_1,\ldots,o_{n_x})}s'$ such that $s$ satisfies the conditions \meta{C}, the objects and registers satisfy the extended conditions \meta{C_\times} in $s$, no object $o_i$ re-enters the same concept, and the objects and registers satisfy the extended effects \meta{E_\times} across $(s,s')$. Registers are unchanged. At least one argument must enter some concept.\\
%Sketch & \code{(:sketch (:conditions \meta{C}) (:effects \meta{E}))} & If $E$ is nonempty, branches over every immediate transition $s\xrightarrow{a}s'$ satisfying $E$. If $E$ is empty, performs one memory-only step with $s'=s$. Registers are unchanged. \\
%Do & \code{(:do (:conditions \meta{C}) (:action "\meta{a}") (:arguments \meta{x_1} \(\ldots\) \meta{x_{n_x}}) (:effects \meta{E}))} & Branches over every ground instance of schema $a$ whose objects belong to the listed concept-feature denotations and whose immediate transition satisfies $E$. \\
Load concept & \code{(:load (:conditions \meta{C}) (:concept \meta{f_C}) (:register (:concept \meta{r})))} & For every $d\in (f_C)_P(s,\rho)$, sets $\rho_C(r)=d$ without changing the concrete state. \\
Load role & \code{(:load (:conditions \meta{C}) (:role \meta{f_R}) (:register (:role \meta{o})))} & For every $(d_1,d_2)\in (f_R)_P(s,\rho)$, sets $\rho_R(o)=(d_1,d_2)$ without changing the concrete state. \\
\bottomrule
\end{longtable}
\end{compacttable}

\section{Execution}

\subsection{Configurations}

An execution configuration is
\[
  \xi=(m,q,\rho,h,s),
\]
where $m$ is the sole module, $q$ is the active memory state, $\rho$ maps loaded
registers to typed values, $h$ is a history of for each concept, and $s\in S_P$ is the concrete planning state.
The history of a concept is a subset of objects.
Two configurations are distinct if any component differs.
Initially execution starts in
$(m,q_0,\emptyset,h_0,s_0)$, where $q_0$ is the module's entry memory, and $h_0$ is the empty history for each concept.

A rule is enabled only when its source memory is current. Load and Concept 
bodies enter its target memory and determine the remaining state and register
effects.

\subsection{Operational Semantics}

A goal state terminates immediately. Otherwise, execution takes the union of
the successors produced by every applicable Load and Concept body of every
enabled rule. No body kind or rule suppresses successors produced by another.
If this full union is empty, execution restores the caller when one exists;
without a caller, this is top-level failure at the non-goal state.

Within this union, universal \code{find_solution} branches over every
applicable body, load value, matching ground action, module call, and immediate concept-rule
successor. Existential \code{find_solution} samples one representative successor
at each choice, so seeds may produce different rollouts. Universal traversal is
the correctness semantics over all choices; existential traversal only obtains
one successful witness or one counterexample and cannot establish correctness.
Both modes consider only immediate successors; neither searches for a multi-step
endpoint satisfying a body. An enabled Load whose referenced feature is empty
is a failing branch.
%An enabled Do with no matching immediate action contributes
%no branch; it neither fails independently nor suppresses branches from other
%applicable rules.
Caller restoration or the no-applicable-action fallback occurs
only when the full successor union is empty.

A load body must target a register declared in the enclosing module and of the
same type as the referenced feature. The selected value becomes available through
\code{(c_register r)} or \code{(r_register o)} in later expressions of the same
module execution.

The history of concepts are augmented with the objects that enter the concept
across the taken transition. When a concept changes, the history of every concept
that precedes it is reset (i.e., set to the empty set).
Transitions across which an object re-enters the same concept are disabled.
The relation defined by the precedence pairs in the reset specification
must be a strict partial order; in particular, acyclic.

\section{Termination}

Reaching a planning goal preempts every body and successfully terminates the
whole program. Away from a goal, an empty full successor union restores the
caller when one exists and is top-level failure otherwise. A maximal program
trajectory is a finite or infinite sequence of configurations generated by
enabled rule bodies, immediate concrete transitions, loads, calls, and caller
restorations.

Runir proves structural termination of the sole module on a finite abstract policy
graph whose vertices combine memory states with qualitative boolean and
numerical feature valuations. Within each strongly connected component, the
Sieve algorithm repeatedly removes an edge justified by an unopposed strict
numerical increase or decrease, or an unopposed boolean change. A module is
structurally terminating when no cycle survives.

The public analysis defaults are \code{max_features=10} and
\code{use_incomplete_preprocessing=True}. Incomplete preprocessing cheaply
removes rules already justified by local qualitative changes before the full
Sieve check. Exceeding \code{max_features} in a residual memory component is an
analysis error, not a nontermination witness.

Rule effects determine the abstract changes used by Sieve.
%Sketch and Do bodies use their declared effects; unmentioned features are unconstrained.
A Load preserves features that do not reference its target register and leaves
register-dependent features unconstrained.

A module program passes the structural check iff its sole module passes. A
reported counterexample is a surviving cycle in the abstract termination graph,
not a concrete execution trace; it identifies the memory states, rules, and
feature changes that still require a ranking argument or structural redesign.

\section{Correctness}

A module program $\Pi$ solves a task $P$ iff every maximal trajectory from
the initial configuration is finite and terminates successfully in a state
$s\in S_G$. Therefore correctness is universal over all nondeterministic choices
left by sketches, loads, actions with the same denoted action arguments, and
applicable rules. A module program solves $\mathcal Q$ iff it solves every solvable
$P\in\mathcal Q$.

\section*{Source}

This document follows the generalized-planning qualitative feature-rule model of
arXiv:2101.00692 and extends it with Runir module-program control states,
registers, and concrete action bodies.

\end{document}
